\Chapter{103}

\Verse{1}
{
    \zR{话说贾琏到了王夫人那边，一一的说了。}
}
{
    \eR{[English source not present in the checked-in Bencraft/Joly files.]}
}
{
}

\Verse{2}
{
    \zR{次日，到了部里，打点停妥，回来又 到王夫人那边将打点吏部之事告知王夫人。}
}
{
    \eR{[English source not present in the checked-in Bencraft/Joly files.]}
}
{
}

\Verse{3}
{
    \zR{王夫人便道：“打听准了么？}
}
{
    \eR{[English source not present in the checked-in Bencraft/Joly files.]}
}
{
}

\Verse{4}
{
    \zR{果然这样， 老爷也愿意，合家也放心。}
}
{
    \eR{[English source not present in the checked-in Bencraft/Joly files.]}
}
{
}

\Verse{5}
{
    \zR{那外任何尝是做得的？}
}
{
    \eR{[English source not present in the checked-in Bencraft/Joly files.]}
}
{
}

\Verse{6}
{
    \zR{不是这样回来，只怕叫那些混帐 东西把老爷的性命都坑了呢。”}
}
{
    \eR{[English source not present in the checked-in Bencraft/Joly files.]}
}
{
}

\Verse{7}
{
    \zR{贾琏道：“太太怎么知道？”}
}
{
    \eR{[English source not present in the checked-in Bencraft/Joly files.]}
}
{
}

\Verse{8}
{
    \zR{王夫人道：“自从你二 叔放了外任，并没有一个钱拿回来，把家里的倒掏摸了好些去了。}
}
{
    \eR{[English source not present in the checked-in Bencraft/Joly files.]}
}
{
}

\Verse{9}
{
    \zR{你瞧那些跟老爷 去的人，他男人在外头不多几时，那些小老婆子们都金头银面的妆扮起来了，可不 是在外头瞒着老爷弄钱？}
}
{
    \eR{[English source not present in the checked-in Bencraft/Joly files.]}
}
{
}

\Verse{10}
{
    \zR{你叔叔就由着他们闹去。}
}
{
    \eR{[English source not present in the checked-in Bencraft/Joly files.]}
}
{
}

\Verse{11}
{
    \zR{要弄出事来，不但自己的官做不 成，只怕连祖上的官也要抹掉了呢。”}
}
{
    \eR{[English source not present in the checked-in Bencraft/Joly files.]}
}
{
}

\Verse{12}
{
    \zR{贾琏道：“太太说的很是。}
}
{
    \eR{[English source not present in the checked-in Bencraft/Joly files.]}
}
{
}

\Verse{13}
{
    \zR{方才我听见参了， 吓的了不得，直等打听明白才放心。}
}
{
    \eR{[English source not present in the checked-in Bencraft/Joly files.]}
}
{
}

\Verse{14}
{
    \zR{也愿意老爷做个京官，安安逸逸的做几年，才 保得住一辈子的声名。}
}
{
    \eR{[English source not present in the checked-in Bencraft/Joly files.]}
}
{
}

\Verse{15}
{
    \zR{就是老太太知道了，倒也是放心的。}
}
{
    \eR{[English source not present in the checked-in Bencraft/Joly files.]}
}
{
}

\Verse{16}
{
    \zR{只要太太说的宽缓些。”}
}
{
    \eR{[English source not present in the checked-in Bencraft/Joly files.]}
}
{
}

\Verse{17}
{
    \zR{王夫人道：“我知道，你到底再去打听打听。”}
}
{
    \eR{[English source not present in the checked-in Bencraft/Joly files.]}
}
{
}

\Verse{18}
{
    \zR{贾琏答应了，才要出来，只见薛姨妈家的老婆子慌慌张张的走来，到王夫人里 间屋内，也没说请安，便道：“我们太太叫我来告诉这里的姨太太说：我们家了不
        得了，又闹出事来了！”}
}
{
    \eR{[English source not present in the checked-in Bencraft/Joly files.]}
}
{
}

\Verse{19}
{
    \zR{王夫人听了，便问：“闹出什么事来？”}
}
{
    \eR{[English source not present in the checked-in Bencraft/Joly files.]}
}
{
}

\Verse{20}
{
    \zR{那婆子又说：“了 不得，了不得！”}
}
{
    \eR{[English source not present in the checked-in Bencraft/Joly files.]}
}
{
}

\Verse{21}
{
    \zR{王夫人哼道：“糊涂东西!有紧要事你到底说呀。”}
}
{
    \eR{[English source not present in the checked-in Bencraft/Joly files.]}
}
{
}

\Verse{22}
{
    \zR{婆子便说：“我 们家二爷不在家，一个男人也没有，这件事情出来，怎么办!要求太太打发几位爷 们去料理料理。”}
}
{
    \eR{[English source not present in the checked-in Bencraft/Joly files.]}
}
{
}

\Verse{23}
{
    \zR{王夫人听着不懂，便着急道：“到底要爷们去干什么？”}
}
{
    \eR{[English source not present in the checked-in Bencraft/Joly files.]}
}
{
}

\Verse{24}
{
    \zR{婆子道： “我们大奶奶死了！”}
}
{
    \eR{[English source not present in the checked-in Bencraft/Joly files.]}
}
{
}

\Verse{25}
{
    \zR{王夫人听了，啐道：“呸，那行子女人死就死了罢咧，也值的 大惊小怪的。”}
}
{
    \eR{[English source not present in the checked-in Bencraft/Joly files.]}
}
{
}

\Verse{26}
{
    \zR{婆子道：“不是好好儿死的，是混闹死的。}
}
{
    \eR{[English source not present in the checked-in Bencraft/Joly files.]}
}
{
}

\Verse{27}
{
    \zR{快求太太打发人去办办！”}
}
{
    \eR{[English source not present in the checked-in Bencraft/Joly files.]}
}
{
}

\Verse{28}
{
    \zR{说着就要走。}
}
{
    \eR{[English source not present in the checked-in Bencraft/Joly files.]}
}
{
}

\Verse{29}
{
    \zR{王夫人又生气，又好笑，说：“这老婆子好混账。}
}
{
    \eR{[English source not present in the checked-in Bencraft/Joly files.]}
}
{
}

\Verse{30}
{
    \zR{琏哥儿，倒不如你 去瞧瞧，别理那糊涂东西。”}
}
{
    \eR{[English source not present in the checked-in Bencraft/Joly files.]}
}
{
}

\Verse{31}
{
    \zR{那婆子没听见打发人去，只听见说“别理他”，他便赌 气跑回去了。}
}
{
    \eR{[English source not present in the checked-in Bencraft/Joly files.]}
}
{
}

\Verse{32}
{
    \zR{这里薛姨妈正在着急，再不见来。}
}
{
    \eR{[English source not present in the checked-in Bencraft/Joly files.]}
}
{
}

\Verse{33}
{
    \zR{好容易那婆子来了，便问：“姨太 太打发谁来？”}
}
{
    \eR{[English source not present in the checked-in Bencraft/Joly files.]}
}
{
}

\Verse{34}
{
    \zR{婆子叹说道：“人再别有急难事。}
}
{
    \eR{[English source not present in the checked-in Bencraft/Joly files.]}
}
{
}

\Verse{35}
{
    \zR{什么好亲好眷，看来也不中用。}
}
{
    \eR{[English source not present in the checked-in Bencraft/Joly files.]}
}
{
}

\Verse{36}
{
    \zR{姨太太不但不肯照应我们，倒骂我糊涂。”}
}
{
    \eR{[English source not present in the checked-in Bencraft/Joly files.]}
}
{
}

\Verse{37}
{
    \zR{薛姨妈听了，又气又急道：“姨太太不管， 你姑奶奶怎么说来着？”}
}
{
    \eR{[English source not present in the checked-in Bencraft/Joly files.]}
}
{
}

\Verse{38}
{
    \zR{婆子道：“姨太太既不管，我们家的姑奶奶自然更不管了， 没有去告诉。”}
}
{
    \eR{[English source not present in the checked-in Bencraft/Joly files.]}
}
{
}

\Verse{39}
{
    \zR{薛姨妈啐道：“姨太太是外人，姑娘是我养的，怎么不管？”}
}
{
    \eR{[English source not present in the checked-in Bencraft/Joly files.]}
}
{
}

\Verse{40}
{
    \zR{婆子一 时省悟道：“是啊，这么着我还去。”}
}
{
    \eR{[English source not present in the checked-in Bencraft/Joly files.]}
}
{
}

\Verse{41}
{
    \zR{正说着，只见贾琏来了，给薛姨妈请了安，道了恼，回说：“我婶子知道弟妇 死了，问老婆子再说不明。}
}
{
    \eR{[English source not present in the checked-in Bencraft/Joly files.]}
}
{
}

\Verse{42}
{
    \zR{着急的很，打发我来问个明白，还叫我在这里料理。}
}
{
    \eR{[English source not present in the checked-in Bencraft/Joly files.]}
}
{
}

\Verse{43}
{
    \zR{该 怎么样，姨太太只管说了办去。”}
}
{
    \eR{[English source not present in the checked-in Bencraft/Joly files.]}
}
{
}

\Verse{44}
{
    \zR{薛姨妈本来气的干哭，听见贾琏的话，便赶忙说： “倒叫二爷费心。}
}
{
    \eR{[English source not present in the checked-in Bencraft/Joly files.]}
}
{
}

\Verse{45}
{
    \zR{我说姨太太是待我最好的，都是这老货说不清，几乎误了事。}
}
{
    \eR{[English source not present in the checked-in Bencraft/Joly files.]}
}
{
}

\Verse{46}
{
    \zR{请 二爷坐下，等我慢慢的告诉你。”}
}
{
    \eR{[English source not present in the checked-in Bencraft/Joly files.]}
}
{
}

\Verse{47}
{
    \zR{便道：“不为别的事，为的是媳妇不是好死的。”}
}
{
    \eR{[English source not present in the checked-in Bencraft/Joly files.]}
}
{
}

\Verse{48}
{
    \zR{贾琏道：“想是为兄弟犯事，怨命死的？”}
}
{
    \eR{[English source not present in the checked-in Bencraft/Joly files.]}
}
{
}

\Verse{49}
{
    \zR{薛姨妈道：“若这样倒好了。}
}
{
    \eR{[English source not present in the checked-in Bencraft/Joly files.]}
}
{
}

\Verse{50}
{
    \zR{前几个月头 里，他天天赤脚蓬头的疯闹。}
}
{
    \eR{[English source not present in the checked-in Bencraft/Joly files.]}
}
{
}

\Verse{51}
{
    \zR{后来听见你兄弟问了死罪，他虽哭了一场，以后倒擦 胭抹粉的起来。}
}
{
    \eR{[English source not present in the checked-in Bencraft/Joly files.]}
}
{
}

\Verse{52}
{
    \zR{我要说他，又要吵个了不得，我总不理他。}
}
{
    \eR{[English source not present in the checked-in Bencraft/Joly files.]}
}
{
}

\Verse{53}
{
    \zR{有一天，不知为什么来 要香菱去作伴儿。}
}
{
    \eR{[English source not present in the checked-in Bencraft/Joly files.]}
}
{
}

\Verse{54}
{
    \zR{我说：‘你放着宝蟾，要香菱做什么？}
}
{
    \eR{[English source not present in the checked-in Bencraft/Joly files.]}
}
{
}

\Verse{55}
{
    \zR{况且香菱是你不爱的，何苦 惹气呢？’}
}
{
    \eR{[English source not present in the checked-in Bencraft/Joly files.]}
}
{
}

\Verse{56}
{
    \zR{他必不依。}
}
{
    \eR{[English source not present in the checked-in Bencraft/Joly files.]}
}
{
}

\Verse{57}
{
    \zR{我没法儿，只得叫香菱到他屋里去。}
}
{
    \eR{[English source not present in the checked-in Bencraft/Joly files.]}
}
{
}

\Verse{58}
{
    \zR{可怜香菱不敢违我的话， 带着病就去了。}
}
{
    \eR{[English source not present in the checked-in Bencraft/Joly files.]}
}
{
}

\Verse{59}
{
    \zR{谁知道他待香菱很好。}
}
{
    \eR{[English source not present in the checked-in Bencraft/Joly files.]}
}
{
}

\Verse{60}
{
    \zR{我倒喜欢，你大妹妹知道了说：‘只怕不是 好心罢？’}
}
{
    \eR{[English source not present in the checked-in Bencraft/Joly files.]}
}
{
}

\Verse{61}
{
    \zR{我也不理会。}
}
{
    \eR{[English source not present in the checked-in Bencraft/Joly files.]}
}
{
}

\Verse{62}
{
    \zR{头几天香菱病着，他倒亲手去做汤给他喝。}
}
{
    \eR{[English source not present in the checked-in Bencraft/Joly files.]}
}
{
}

\Verse{63}
{
    \zR{谁知香菱没福， 刚端到跟前，他自己烫了手，连碗都砸了。}
}
{
    \eR{[English source not present in the checked-in Bencraft/Joly files.]}
}
{
}

\Verse{64}
{
    \zR{我只说必要迁怒在香菱身上，他倒没生 气，自己还拿笤帚扫了，拿水泼净了地，仍旧两个人很好。}
}
{
    \eR{[English source not present in the checked-in Bencraft/Joly files.]}
}
{
}

\Verse{65}
{
    \zR{昨儿晚上，又叫宝蟾去 做了两碗汤来，自己说和香菱一块儿喝。}
}
{
    \eR{[English source not present in the checked-in Bencraft/Joly files.]}
}
{
}

\Verse{66}
{
    \zR{隔了一会子，听见他屋里闹起来，宝蟾急 的乱嚷，以后香菱也嚷着，扶着墙出来叫人。}
}
{
    \eR{[English source not present in the checked-in Bencraft/Joly files.]}
}
{
}

\Verse{67}
{
    \zR{我忙着看去，只见媳妇鼻子眼睛里都 流出血来，在地下乱滚，两只手在心口里乱抓，两只脚乱蹬，把我就吓死了。}
}
{
    \eR{[English source not present in the checked-in Bencraft/Joly files.]}
}
{
}

\Verse{68}
{
    \zR{问他 也说不出来，闹了一会子就死了。}
}
{
    \eR{[English source not present in the checked-in Bencraft/Joly files.]}
}
{
}

\Verse{69}
{
    \zR{我瞧那个光景儿是服了毒的。}
}
{
    \eR{[English source not present in the checked-in Bencraft/Joly files.]}
}
{
}

\Verse{70}
{
    \zR{宝蟾就哭着来揪香 菱，说他拿药药死奶奶了。}
}
{
    \eR{[English source not present in the checked-in Bencraft/Joly files.]}
}
{
}

\Verse{71}
{
    \zR{我看香菱也不是这么样的人，再者他病的起还起不来， 怎么能药人呢？}
}
{
    \eR{[English source not present in the checked-in Bencraft/Joly files.]}
}
{
}

\Verse{72}
{
    \zR{无奈宝蟾一口咬定，我的二爷，这叫我怎么办？}
}
{
    \eR{[English source not present in the checked-in Bencraft/Joly files.]}
}
{
}

\Verse{73}
{
    \zR{只得硬着心肠叫老婆 子们把香菱捆了，交给宝蟾，便把房门反扣了。}
}
{
    \eR{[English source not present in the checked-in Bencraft/Joly files.]}
}
{
}

\Verse{74}
{
    \zR{我和你二妹妹守了一夜，等府里的 门开了才告诉去的。}
}
{
    \eR{[English source not present in the checked-in Bencraft/Joly files.]}
}
{
}

\Verse{75}
{
    \zR{二爷你是明白人，这件事怎么好？”}
}
{
    \eR{[English source not present in the checked-in Bencraft/Joly files.]}
}
{
}

\Verse{76}
{
    \zR{贾琏道：“夏家知道了没 有？”}
}
{
    \eR{[English source not present in the checked-in Bencraft/Joly files.]}
}
{
}

\Verse{77}
{
    \zR{薛姨妈道：“也得撕掳明白了，才好报啊。”}
}
{
    \eR{[English source not present in the checked-in Bencraft/Joly files.]}
}
{
}

\Verse{78}
{
    \zR{贾琏道：“据我看起来，必要经 官才了的下来。}
}
{
    \eR{[English source not present in the checked-in Bencraft/Joly files.]}
}
{
}

\Verse{79}
{
    \zR{我们自然疑在宝蟾身上，别人却说宝蟾为什么药死他们姑娘呢？}
}
{
    \eR{[English source not present in the checked-in Bencraft/Joly files.]}
}
{
}

\Verse{80}
{
    \zR{若 说在香菱身上，倒还装得上。”}
}
{
    \eR{[English source not present in the checked-in Bencraft/Joly files.]}
}
{
}

\Verse{81}
{
    \zR{正说着，只见荣府的女人们进来说：“我们二奶奶来了。”}
}
{
    \eR{[English source not present in the checked-in Bencraft/Joly files.]}
}
{
}

\Verse{82}
{
    \zR{贾琏虽是大伯子，因 从小儿见的，也不回避。}
}
{
    \eR{[English source not present in the checked-in Bencraft/Joly files.]}
}
{
}

\Verse{83}
{
    \zR{宝钗进来见了母亲，又见了贾琏，便往里间屋里和宝琴坐 下。}
}
{
    \eR{[English source not present in the checked-in Bencraft/Joly files.]}
}
{
}

\Verse{84}
{
    \zR{薛姨妈进来也将前事告诉了一遍。}
}
{
    \eR{[English source not present in the checked-in Bencraft/Joly files.]}
}
{
}

\Verse{85}
{
    \zR{宝钗便说：“若把香菱捆了，可不是我们也 说是香菱药死的了么？}
}
{
    \eR{[English source not present in the checked-in Bencraft/Joly files.]}
}
{
}

\Verse{86}
{
    \zR{妈妈说这汤是宝蟾做的，就该捆起宝蟾来问他呀。}
}
{
    \eR{[English source not present in the checked-in Bencraft/Joly files.]}
}
{
}

\Verse{87}
{
    \zR{一面就该 打发人报夏家去，一面报官才是。”}
}
{
    \eR{[English source not present in the checked-in Bencraft/Joly files.]}
}
{
}

\Verse{88}
{
    \zR{薛姨妈听见有理，便问贾琏。}
}
{
    \eR{[English source not present in the checked-in Bencraft/Joly files.]}
}
{
}

\Verse{89}
{
    \zR{贾琏道：“二妹子 说的很是。}
}
{
    \eR{[English source not present in the checked-in Bencraft/Joly files.]}
}
{
}

\Verse{90}
{
    \zR{报官还得我去托了刑部里的人，相验问口供的时候，方有照应。}
}
{
    \eR{[English source not present in the checked-in Bencraft/Joly files.]}
}
{
}

\Verse{91}
{
    \zR{只是要 捆宝蟾放香菱，倒怕难些。”}
}
{
    \eR{[English source not present in the checked-in Bencraft/Joly files.]}
}
{
}

\Verse{92}
{
    \zR{薛姨妈道：“并不是我要捆香菱，我恐怕香菱病中受冤 着急，一时寻死，又添了一条人命，才捆了交给宝蟾，也是个主意。”}
}
{
    \eR{[English source not present in the checked-in Bencraft/Joly files.]}
}
{
}

\Verse{93}
{
    \zR{贾琏道：“虽 是这么说，我们倒帮了宝蟾了。}
}
{
    \eR{[English source not present in the checked-in Bencraft/Joly files.]}
}
{
}

\Verse{94}
{
    \zR{若要放都放，要捆都捆，他们三个人是一处的。}
}
{
    \eR{[English source not present in the checked-in Bencraft/Joly files.]}
}
{
}

\Verse{95}
{
    \zR{只 要叫人安慰香菱就是了。”}
}
{
    \eR{[English source not present in the checked-in Bencraft/Joly files.]}
}
{
}

\Verse{96}
{
    \zR{薛姨妈便叫人开门进去。}
}
{
    \eR{[English source not present in the checked-in Bencraft/Joly files.]}
}
{
}

\Verse{97}
{
    \zR{宝钗就派了带来的几个女人帮 着捆宝蟾。}
}
{
    \eR{[English source not present in the checked-in Bencraft/Joly files.]}
}
{
}

\Verse{98}
{
    \zR{只见香菱已哭的死去活来。}
}
{
    \eR{[English source not present in the checked-in Bencraft/Joly files.]}
}
{
}

\Verse{99}
{
    \zR{宝蟾反得意洋洋，以后见人要捆他，便乱嚷 起来，那禁得荣府的人吆喝着，也就捆了，竟开着门，好叫人看着。}
}
{
    \eR{[English source not present in the checked-in Bencraft/Joly files.]}
}
{
}

\Verse{100}
{
    \zR{这里报夏家的 人已经去了。}
}
{
    \eR{[English source not present in the checked-in Bencraft/Joly files.]}
}
{
}

\Verse{101}
{
    \zR{那夏家先前不住在京里，因近年消索，又惦记女孩儿，新近搬进京来。}
}
{
    \eR{[English source not present in the checked-in Bencraft/Joly files.]}
}
{
}

\Verse{102}
{
    \zR{父亲已 没，只有母亲，又过继了一个混账儿子，把家业都花完了，不时的常到薛家。}
}
{
    \eR{[English source not present in the checked-in Bencraft/Joly files.]}
}
{
}

\Verse{103}
{
    \zR{那金 桂原是个水性人儿，那里守得住空房，况兼天天心里想念薛蝌，便有些饥不择食的 光景。}
}
{
    \eR{[English source not present in the checked-in Bencraft/Joly files.]}
}
{
}

\Verse{104}
{
    \zR{无奈他这个干兄弟又是个蠢货，虽也有些知觉，只是尚未入港，所以金桂时 常回去，也帮贴他些银钱。}
}
{
    \eR{[English source not present in the checked-in Bencraft/Joly files.]}
}
{
}

\Verse{105}
{
    \zR{这些时正盼金桂回家，只见薛家的人来，心里想着：“又 拿什么东西来了。”}
}
{
    \eR{[English source not present in the checked-in Bencraft/Joly files.]}
}
{
}

\Verse{106}
{
    \zR{不料说这里的姑娘服毒死了，他就气的乱嚷乱叫。}
}
{
    \eR{[English source not present in the checked-in Bencraft/Joly files.]}
}
{
}

\Verse{107}
{
    \zR{金桂的母亲 听见了，更哭喊起来，说：“好端端的女孩儿在他家，为什么服了毒呢！”}
}
{
    \eR{[English source not present in the checked-in Bencraft/Joly files.]}
}
{
}

\Verse{108}
{
    \zR{哭着喊着 的，带了儿子，也等不得雇车，便要走来。}
}
{
    \eR{[English source not present in the checked-in Bencraft/Joly files.]}
}
{
}

\Verse{109}
{
    \zR{那夏家本是买卖人家，如今没了钱，那 顾什么脸面，儿子头里走，他就跟了个破老婆子出了门，在街上哭哭啼啼的雇了一 辆车，一直跑到薛家。}
}
{
    \eR{[English source not present in the checked-in Bencraft/Joly files.]}
}
{
}

\Verse{110}
{
    \zR{进门也不搭话，就“儿”一声“肉”一声的闹起。}
}
{
    \eR{[English source not present in the checked-in Bencraft/Joly files.]}
}
{
}

\Verse{111}
{
    \zR{那时贾琏 到刑部去托人，家里只有薛姨妈、宝钗、宝琴，何曾见过这个阵仗儿，都吓的不敢 则声。}
}
{
    \eR{[English source not present in the checked-in Bencraft/Joly files.]}
}
{
}

\Verse{112}
{
    \zR{要和他讲理，他也不听，只说：“我女孩儿在你家，得过什么好处？}
}
{
    \eR{[English source not present in the checked-in Bencraft/Joly files.]}
}
{
}

\Verse{113}
{
    \zR{两口子朝 打暮骂，闹了几时，还不容他两口子在一处。}
}
{
    \eR{[English source not present in the checked-in Bencraft/Joly files.]}
}
{
}

\Verse{114}
{
    \zR{你们商量着把我女婿弄在监里，永不 见面。}
}
{
    \eR{[English source not present in the checked-in Bencraft/Joly files.]}
}
{
}

\Verse{115}
{
    \zR{你们娘儿们仗着好亲戚受用也罢了，还嫌他碍眼，叫人药死他，倒说是服毒! 他为什么服毒？”}
}
{
    \eR{[English source not present in the checked-in Bencraft/Joly files.]}
}
{
}

\Verse{116}
{
    \zR{说着，直奔薛姨妈来。}
}
{
    \eR{[English source not present in the checked-in Bencraft/Joly files.]}
}
{
}

\Verse{117}
{
    \zR{薛姨妈只得退后，说：“亲家太太!且瞧瞧 你女孩儿，问问宝蟾，再说歪话还不迟呢！”}
}
{
    \eR{[English source not present in the checked-in Bencraft/Joly files.]}
}
{
}

\Verse{118}
{
    \zR{宝钗宝琴因外面有夏家的儿子，难以 出来拦护，只在里边着急。}
}
{
    \eR{[English source not present in the checked-in Bencraft/Joly files.]}
}
{
}

\Verse{119}
{
    \zR{恰好王夫人打发周瑞家的照看，一进门来，见一个老婆子指着薛姨妈的脸哭骂。}
}
{
    \eR{[English source not present in the checked-in Bencraft/Joly files.]}
}
{
}

\Verse{120}
{
    \zR{周瑞家的知道必是金桂的母亲，便走上来说：“这位是亲家太太么？}
}
{
    \eR{[English source not present in the checked-in Bencraft/Joly files.]}
}
{
}

\Verse{121}
{
    \zR{大奶奶自己服毒 死的，与我们姨太太什么相干？}
}
{
    \eR{[English source not present in the checked-in Bencraft/Joly files.]}
}
{
}

\Verse{122}
{
    \zR{也不犯这么遭塌呀。”}
}
{
    \eR{[English source not present in the checked-in Bencraft/Joly files.]}
}
{
}

\Verse{123}
{
    \zR{那金桂的母亲问：“你是谁？”}
}
{
    \eR{[English source not present in the checked-in Bencraft/Joly files.]}
}
{
}

\Verse{124}
{
    \zR{薛姨妈见有了人，胆子略壮了些，便说：“这就是我们亲戚贾府里的。”}
}
{
    \eR{[English source not present in the checked-in Bencraft/Joly files.]}
}
{
}

\Verse{125}
{
    \zR{金桂的母亲 便道：“谁不知道你们有仗腰子的亲戚，才能够叫姑爷坐在监里!如今我的女孩儿倒 白死了不成？”}
}
{
    \eR{[English source not present in the checked-in Bencraft/Joly files.]}
}
{
}

\Verse{126}
{
    \zR{说着，便拉薛姨妈说：“你到底把我女孩儿怎么弄杀了？}
}
{
    \eR{[English source not present in the checked-in Bencraft/Joly files.]}
}
{
}

\Verse{127}
{
    \zR{给我瞧瞧！”}
}
{
    \eR{[English source not present in the checked-in Bencraft/Joly files.]}
}
{
}

\Verse{128}
{
    \zR{周瑞家的一面劝说：“只管瞧去，不用拉拉扯扯。”}
}
{
    \eR{[English source not present in the checked-in Bencraft/Joly files.]}
}
{
}

\Verse{129}
{
    \zR{把手只一推。}
}
{
    \eR{[English source not present in the checked-in Bencraft/Joly files.]}
}
{
}

\Verse{130}
{
    \zR{夏家的儿子便跑进 来不依，道：“你仗着府里的势头儿来打我母亲么？”}
}
{
    \eR{[English source not present in the checked-in Bencraft/Joly files.]}
}
{
}

\Verse{131}
{
    \zR{说着，便将椅子打去，却没 有打着。}
}
{
    \eR{[English source not present in the checked-in Bencraft/Joly files.]}
}
{
}

\Verse{132}
{
    \zR{里头跟宝钗的人听见外头闹起来，赶着来瞧，恐怕周瑞家的吃亏，齐打伙 儿上去，半劝半喝。}
}
{
    \eR{[English source not present in the checked-in Bencraft/Joly files.]}
}
{
}

\Verse{133}
{
    \zR{那夏家的母子，索性撒起泼来，说：“知道你们荣府的势头儿! 我们家的姑娘已经死了，如今也都不要命了！”}
}
{
    \eR{[English source not present in the checked-in Bencraft/Joly files.]}
}
{
}

\Verse{134}
{
    \zR{说着，仍奔薛姨妈拚命。}
}
{
    \eR{[English source not present in the checked-in Bencraft/Joly files.]}
}
{
}

\Verse{135}
{
    \zR{地下的人 虽多，那里挡得住，自古说的：“一人拚命，万夫莫当。”}
}
{
    \eR{[English source not present in the checked-in Bencraft/Joly files.]}
}
{
}

\Verse{136}
{
    \zR{正闹到危急之际，贾琏带了七八个家人进来，见是如此，便叫人先把夏家的儿 子拉出去，便说：“你们不许闹，有话好好儿的说。}
}
{
    \eR{[English source not present in the checked-in Bencraft/Joly files.]}
}
{
}

\Verse{137}
{
    \zR{快将家里收拾收拾，刑部里头 的老爷们就来相验了。”}
}
{
    \eR{[English source not present in the checked-in Bencraft/Joly files.]}
}
{
}

\Verse{138}
{
    \zR{金桂的母亲正在撒泼，只见来了一位老爷，几个在头里吆 喝，那些人都垂手侍立。}
}
{
    \eR{[English source not present in the checked-in Bencraft/Joly files.]}
}
{
}

\Verse{139}
{
    \zR{金桂的母亲见这个光景，也不知是贾府何人。}
}
{
    \eR{[English source not present in the checked-in Bencraft/Joly files.]}
}
{
}

\Verse{140}
{
    \zR{又见他儿子 已被众人揪住，又听见说刑部来验，他心里原想看见女孩儿的尸首，先闹个稀烂， 再去喊冤，不承望这里先报了官，也便软了些。}
}
{
    \eR{[English source not present in the checked-in Bencraft/Joly files.]}
}
{
}

\Verse{141}
{
    \zR{薛姨妈已吓糊涂了，还是周瑞家的 回说：“他们来了也没去瞧瞧他们姑娘，便作践起姨太太来了。}
}
{
    \eR{[English source not present in the checked-in Bencraft/Joly files.]}
}
{
}

\Verse{142}
{
    \zR{我们为好劝他，那 里跑进一个野男人，在奶奶们里头混撒村混打，这可不是没有王法了！”}
}
{
    \eR{[English source not present in the checked-in Bencraft/Joly files.]}
}
{
}

\Verse{143}
{
    \zR{贾琏道：“这 会子不用和他讲理，等回来打着问他，说：男人有男人的地方儿，里头都是些姑娘 奶奶们。}
}
{
    \eR{[English source not present in the checked-in Bencraft/Joly files.]}
}
{
}

\Verse{144}
{
    \zR{况且有他母亲还瞧不见他们姑娘么？}
}
{
    \eR{[English source not present in the checked-in Bencraft/Joly files.]}
}
{
}

\Verse{145}
{
    \zR{他跑进来不是要打抢来了么！”}
}
{
    \eR{[English source not present in the checked-in Bencraft/Joly files.]}
}
{
}

\Verse{146}
{
    \zR{家人们 做好做歹，压伏住了。}
}
{
    \eR{[English source not present in the checked-in Bencraft/Joly files.]}
}
{
}

\Verse{147}
{
    \zR{周瑞家的仗着人多，便说：“夏太太，你不懂事!既来了，该 问个青红皂白。}
}
{
    \eR{[English source not present in the checked-in Bencraft/Joly files.]}
}
{
}

\Verse{148}
{
    \zR{你们姑娘是自己服毒死了，不然就是宝蟾药死他主子了。}
}
{
    \eR{[English source not present in the checked-in Bencraft/Joly files.]}
}
{
}

\Verse{149}
{
    \zR{怎么不问 明白，又不看尸首，就想讹人来了呢？}
}
{
    \eR{[English source not present in the checked-in Bencraft/Joly files.]}
}
{
}

\Verse{150}
{
    \zR{我们就肯叫一个媳妇儿白死了不成？}
}
{
    \eR{[English source not present in the checked-in Bencraft/Joly files.]}
}
{
}

\Verse{151}
{
    \zR{现在把宝 蟾捆着，因为你们姑娘必要点病儿，所以叫香菱陪着他，也在一个屋里住，故此两 个人都看守在那里。}
}
{
    \eR{[English source not present in the checked-in Bencraft/Joly files.]}
}
{
}

\Verse{152}
{
    \zR{原等你们来眼看着刑部相验，问出道理来才是啊。”}
}
{
    \eR{[English source not present in the checked-in Bencraft/Joly files.]}
}
{
}

\Verse{153}
{
    \zR{金桂的母 亲此时势孤，也只得跟着周瑞家的到他女孩儿屋里，只见满脸黑血，直挺挺的躺在 炕上，便叫哭起来。}
}
{
    \eR{[English source not present in the checked-in Bencraft/Joly files.]}
}
{
}

\Verse{154}
{
    \zR{宝蟾见是他家的人来，便哭喊说：“我们姑娘好意待香菱，叫 他在一块儿住，他倒抽空儿药死我们姑娘！”}
}
{
    \eR{[English source not present in the checked-in Bencraft/Joly files.]}
}
{
}

\Verse{155}
{
    \zR{那时薛家上下人等俱在，便齐声吆喝 道：“胡说!昨日奶奶喝了汤才药死的，这汤可不是你做的？”}
}
{
    \eR{[English source not present in the checked-in Bencraft/Joly files.]}
}
{
}

\Verse{156}
{
    \zR{宝蟾道：“汤是我做 的，端了来，我有事走了。}
}
{
    \eR{[English source not present in the checked-in Bencraft/Joly files.]}
}
{
}

\Verse{157}
{
    \zR{不知香菱起来放了些什么在里头，药死的。”}
}
{
    \eR{[English source not present in the checked-in Bencraft/Joly files.]}
}
{
}

\Verse{158}
{
    \zR{金桂的母 亲没听完，就奔香菱，众人拦住。}
}
{
    \eR{[English source not present in the checked-in Bencraft/Joly files.]}
}
{
}

\Verse{159}
{
    \zR{薛姨妈便道：“这样子是砒霜药的，家里决无此 物。}
}
{
    \eR{[English source not present in the checked-in Bencraft/Joly files.]}
}
{
}

\Verse{160}
{
    \zR{不管香菱宝蟾，终有替他买的，回来刑部少不得问出来，才赖不去。}
}
{
    \eR{[English source not present in the checked-in Bencraft/Joly files.]}
}
{
}

\Verse{161}
{
    \zR{如今把媳 妇权放平正，好等官来相验。”}
}
{
    \eR{[English source not present in the checked-in Bencraft/Joly files.]}
}
{
}

\Verse{162}
{
    \zR{众婆子上来抬放。}
}
{
    \eR{[English source not present in the checked-in Bencraft/Joly files.]}
}
{
}

\Verse{163}
{
    \zR{宝钗道：“都是男人进来，你们将 女人动用的东西检点检点。”}
}
{
    \eR{[English source not present in the checked-in Bencraft/Joly files.]}
}
{
}

\Verse{164}
{
    \zR{只见炕褥底下有一个揉成团的纸包儿。}
}
{
    \eR{[English source not present in the checked-in Bencraft/Joly files.]}
}
{
}

\Verse{165}
{
    \zR{金桂的母亲瞧 见，便拾起打开看时，并没有什么，便撩开了。}
}
{
    \eR{[English source not present in the checked-in Bencraft/Joly files.]}
}
{
}

\Verse{166}
{
    \zR{宝蟾看见道：“可不是有了凭据了! 这个纸包儿我认得：头几天耗子闹的慌，奶奶家去找舅爷要的，拿回来搁在首饰匣 内。}
}
{
    \eR{[English source not present in the checked-in Bencraft/Joly files.]}
}
{
}

\Verse{167}
{
    \zR{必是香菱看见了，拿来药死奶奶的。}
}
{
    \eR{[English source not present in the checked-in Bencraft/Joly files.]}
}
{
}

\Verse{168}
{
    \zR{若不信，你们看看首饰匣里有没有了。”}
}
{
    \eR{[English source not present in the checked-in Bencraft/Joly files.]}
}
{
}

\Verse{169}
{
    \zR{金桂的母亲便依着宝蟾的话，取出匣子来，只有几支银簪子。}
}
{
    \eR{[English source not present in the checked-in Bencraft/Joly files.]}
}
{
}

\Verse{170}
{
    \zR{薛姨妈便说：“怎 么好些首饰都没有了？”}
}
{
    \eR{[English source not present in the checked-in Bencraft/Joly files.]}
}
{
}

\Verse{171}
{
    \zR{宝钗叫人打开箱柜，俱是空的，便道：“嫂子这些东西被 谁拿去？}
}
{
    \eR{[English source not present in the checked-in Bencraft/Joly files.]}
}
{
}

\Verse{172}
{
    \zR{这可要问宝蟾。”}
}
{
    \eR{[English source not present in the checked-in Bencraft/Joly files.]}
}
{
}

\Verse{173}
{
    \zR{金桂的母亲心里也虚了好些，见薛姨妈查问宝蟾，便说： “姑娘的东西，他那里知道？”}
}
{
    \eR{[English source not present in the checked-in Bencraft/Joly files.]}
}
{
}

\Verse{174}
{
    \zR{周瑞家的道：“亲家太太别这么说么。}
}
{
    \eR{[English source not present in the checked-in Bencraft/Joly files.]}
}
{
}

\Verse{175}
{
    \zR{我知道宝姑 娘是天天跟着大奶奶的，怎么说不知道？”}
}
{
    \eR{[English source not present in the checked-in Bencraft/Joly files.]}
}
{
}

\Verse{176}
{
    \zR{宝蟾见问得紧，又不好胡赖，只得说道： “奶奶自己每每带回家去，我管得么？”}
}
{
    \eR{[English source not present in the checked-in Bencraft/Joly files.]}
}
{
}

\Verse{177}
{
    \zR{众人便说：“好个亲家太太!哄着拿姑娘的 东西，哄完了叫他寻死来讹我们。}
}
{
    \eR{[English source not present in the checked-in Bencraft/Joly files.]}
}
{
}

\Verse{178}
{
    \zR{好罢咧，回来相验，就是这么说。”}
}
{
    \eR{[English source not present in the checked-in Bencraft/Joly files.]}
}
{
}

\Verse{179}
{
    \zR{宝钗叫人：“到 外头告诉琏二爷说：别放了夏家的人。”}
}
{
    \eR{[English source not present in the checked-in Bencraft/Joly files.]}
}
{
}

\Verse{180}
{
    \zR{里头金桂的母亲忙了手脚，便骂宝蟾道：“小 蹄子，别嚼舌头了!姑娘几时拿东西到我家去？”}
}
{
    \eR{[English source not present in the checked-in Bencraft/Joly files.]}
}
{
}

\Verse{181}
{
    \zR{宝蟾道：“如今东西是小，给姑娘 偿命是大。”}
}
{
    \eR{[English source not present in the checked-in Bencraft/Joly files.]}
}
{
}

\Verse{182}
{
    \zR{宝琴道：“有了东西，就有偿命的人了。}
}
{
    \eR{[English source not present in the checked-in Bencraft/Joly files.]}
}
{
}

\Verse{183}
{
    \zR{快请琏二哥哥问准了夏家的儿 子买砒霜的话，回来好回刑部里的话。”}
}
{
    \eR{[English source not present in the checked-in Bencraft/Joly files.]}
}
{
}

\Verse{184}
{
    \zR{金桂的母亲着了急道：“这宝蟾必是撞见鬼 了，混说起来。}
}
{
    \eR{[English source not present in the checked-in Bencraft/Joly files.]}
}
{
}

\Verse{185}
{
    \zR{我们姑娘何尝买过砒霜？}
}
{
    \eR{[English source not present in the checked-in Bencraft/Joly files.]}
}
{
}

\Verse{186}
{
    \zR{要这么说，必是宝蟾药死了的！”}
}
{
    \eR{[English source not present in the checked-in Bencraft/Joly files.]}
}
{
}

\Verse{187}
{
    \zR{宝蟾急的 乱嚷，说：“别人赖我也罢了，怎么你们也赖起我来呢？}
}
{
    \eR{[English source not present in the checked-in Bencraft/Joly files.]}
}
{
}

\Verse{188}
{
    \zR{你们不是常和姑娘说，叫他 别受委屈，闹得他们家破人亡，那时将东西卷包儿一走，再配一个好姑爷。}
}
{
    \eR{[English source not present in the checked-in Bencraft/Joly files.]}
}
{
}

\Verse{189}
{
    \zR{这个话 是有的没有？”}
}
{
    \eR{[English source not present in the checked-in Bencraft/Joly files.]}
}
{
}

\Verse{190}
{
    \zR{金桂的母亲还未及答言，周瑞家的便接口说道：“这是你们家的人 说的，还赖什么呢？”}
}
{
    \eR{[English source not present in the checked-in Bencraft/Joly files.]}
}
{
}

\Verse{191}
{
    \zR{金桂的母亲恨的咬牙切齿的骂宝蟾，说：“我待你不错呀， 为什么你倒拿话来葬送我呢？}
}
{
    \eR{[English source not present in the checked-in Bencraft/Joly files.]}
}
{
}

\Verse{192}
{
    \zR{回来见了官，我就说是你药死姑娘的！”}
}
{
    \eR{[English source not present in the checked-in Bencraft/Joly files.]}
}
{
}

\Verse{193}
{
    \zR{宝蟾气的瞪着眼说：“请太太放了香菱罢，不犯着白害别人，我见官自有我的 话。”}
}
{
    \eR{[English source not present in the checked-in Bencraft/Joly files.]}
}
{
}

\Verse{194}
{
    \zR{宝钗听出这个话头儿来了，便叫人反倒放开了宝蟾，说：“你原是个爽快人， 何苦白冤在里头？}
}
{
    \eR{[English source not present in the checked-in Bencraft/Joly files.]}
}
{
}

\Verse{195}
{
    \zR{你有话，索性说了大家明白，岂不完了事了呢？”}
}
{
    \eR{[English source not present in the checked-in Bencraft/Joly files.]}
}
{
}

\Verse{196}
{
    \zR{宝蟾也怕见官 受苦，便说：“我们奶奶天天抱怨说：‘我这样人，为什么碰着这个瞎眼的娘，不配 给二爷，偏给了这么个混账糊涂行子。}
}
{
    \eR{[English source not present in the checked-in Bencraft/Joly files.]}
}
{
}

\Verse{197}
{
    \zR{要是能够和二爷过一天，死了也是愿意的。’}
}
{
    \eR{[English source not present in the checked-in Bencraft/Joly files.]}
}
{
}

\Verse{198}
{
    \zR{说到那里，便恨香菱。}
}
{
    \eR{[English source not present in the checked-in Bencraft/Joly files.]}
}
{
}

\Verse{199}
{
    \zR{我起初不理会，后来看见和香菱好了，我只道是香菱怎么哄 转了。}
}
{
    \eR{[English source not present in the checked-in Bencraft/Joly files.]}
}
{
}

\Verse{200}
{
    \zR{不承望昨儿的汤不是好意。”}
}
{
    \eR{[English source not present in the checked-in Bencraft/Joly files.]}
}
{
}

\Verse{201}
{
    \zR{金桂的母亲接说道：“越发胡说了!若是要药香 菱，为什么倒药了自己呢？”}
}
{
    \eR{[English source not present in the checked-in Bencraft/Joly files.]}
}
{
}

\Verse{202}
{
    \zR{宝钗便问道：“香菱，昨日你喝汤来着没有？”}
}
{
    \eR{[English source not present in the checked-in Bencraft/Joly files.]}
}
{
}

\Verse{203}
{
    \zR{香菱 道：“头几天我病的抬不起头来，奶奶叫我喝汤，我不敢说不喝。}
}
{
    \eR{[English source not present in the checked-in Bencraft/Joly files.]}
}
{
}

\Verse{204}
{
    \zR{刚要扎挣起来， 那碗汤已经洒了，倒叫奶奶收拾了个难，我心里很过不去。}
}
{
    \eR{[English source not present in the checked-in Bencraft/Joly files.]}
}
{
}

\Verse{205}
{
    \zR{昨儿听见叫我喝汤，我 喝不下去，没有法儿，正要喝的时候儿，偏又头晕起来。}
}
{
    \eR{[English source not present in the checked-in Bencraft/Joly files.]}
}
{
}

\Verse{206}
{
    \zR{见宝蟾姐姐端了去。}
}
{
    \eR{[English source not present in the checked-in Bencraft/Joly files.]}
}
{
}

\Verse{207}
{
    \zR{我正 喜欢，刚合上眼，奶奶自己喝着汤，叫我尝尝，我便勉强也喝了两口。”}
}
{
    \eR{[English source not present in the checked-in Bencraft/Joly files.]}
}
{
}

\Verse{208}
{
    \zR{宝蟾不待 说完便道：“是了!我老实说罢。}
}
{
    \eR{[English source not present in the checked-in Bencraft/Joly files.]}
}
{
}

\Verse{209}
{
    \zR{昨儿奶奶叫我做两碗汤，说是和香菱同喝。}
}
{
    \eR{[English source not present in the checked-in Bencraft/Joly files.]}
}
{
}

\Verse{210}
{
    \zR{我气不 过，心里想着：香菱那里配我做汤给他喝呢？}
}
{
    \eR{[English source not present in the checked-in Bencraft/Joly files.]}
}
{
}

\Verse{211}
{
    \zR{我故意的一碗里头多抓了一把盐，记 了暗记儿，原想给香菱喝的。}
}
{
    \eR{[English source not present in the checked-in Bencraft/Joly files.]}
}
{
}

\Verse{212}
{
    \zR{刚端进来，奶奶却拦着我叫外头叫小子们雇车，说今 日回家去。}
}
{
    \eR{[English source not present in the checked-in Bencraft/Joly files.]}
}
{
}

\Verse{213}
{
    \zR{我出去说了回来，见盐多的这碗汤在奶奶跟前呢。}
}
{
    \eR{[English source not present in the checked-in Bencraft/Joly files.]}
}
{
}

\Verse{214}
{
    \zR{我恐怕奶奶喝着咸， 又要骂我。}
}
{
    \eR{[English source not present in the checked-in Bencraft/Joly files.]}
}
{
}

\Verse{215}
{
    \zR{正没法的时候，奶奶往后头走动，我眼错不见，就把香菱这碗汤换过来 了。}
}
{
    \eR{[English source not present in the checked-in Bencraft/Joly files.]}
}
{
}

\Verse{216}
{
    \zR{也是合该如此。}
}
{
    \eR{[English source not present in the checked-in Bencraft/Joly files.]}
}
{
}

\Verse{217}
{
    \zR{奶奶回来就拿了汤去到香菱床边，喝着说：‘你到底尝尝。’}
}
{
    \eR{[English source not present in the checked-in Bencraft/Joly files.]}
}
{
}

\Verse{218}
{
    \zR{那 香菱也不觉咸，两个人都喝完了。}
}
{
    \eR{[English source not present in the checked-in Bencraft/Joly files.]}
}
{
}

\Verse{219}
{
    \zR{我正笑香菱没嘴道儿，那里知道这死鬼奶奶要药 香菱，必定趁我不在，将砒霜撒上了，也不知道我换碗。}
}
{
    \eR{[English source not present in the checked-in Bencraft/Joly files.]}
}
{
}

\Verse{220}
{
    \zR{这可就是天理昭彰，自害 自身了。”}
}
{
    \eR{[English source not present in the checked-in Bencraft/Joly files.]}
}
{
}

\Verse{221}
{
    \zR{于是众人往前后一想，真正一丝不错，便将香菱也放了，扶着他仍旧睡 在床上。}
}
{
    \eR{[English source not present in the checked-in Bencraft/Joly files.]}
}
{
}

\Verse{222}
{
    \zR{不说香菱得放，且说金桂的母亲心虚事实，还想辩赖。}
}
{
    \eR{[English source not present in the checked-in Bencraft/Joly files.]}
}
{
}

\Verse{223}
{
    \zR{薛姨妈等你言我语，反 要他儿子偿还金桂之命。}
}
{
    \eR{[English source not present in the checked-in Bencraft/Joly files.]}
}
{
}

\Verse{224}
{
    \zR{正然吵嚷，贾琏在外嚷说：“不用多说了，快收拾停当。}
}
{
    \eR{[English source not present in the checked-in Bencraft/Joly files.]}
}
{
}

\Verse{225}
{
    \zR{刑部的老爷就到了。”}
}
{
    \eR{[English source not present in the checked-in Bencraft/Joly files.]}
}
{
}

\Verse{226}
{
    \zR{此时惟有夏家母子着忙，想来总要吃亏的，不得已反求薛姨 妈道：“千不是，万不是，总是我死的女孩儿不长进。}
}
{
    \eR{[English source not present in the checked-in Bencraft/Joly files.]}
}
{
}

\Verse{227}
{
    \zR{这也是他自作自受。}
}
{
    \eR{[English source not present in the checked-in Bencraft/Joly files.]}
}
{
}

\Verse{228}
{
    \zR{要是刑 部相验，到底府上脸面不好看，求亲家太太息了这件事罢。”}
}
{
    \eR{[English source not present in the checked-in Bencraft/Joly files.]}
}
{
}

\Verse{229}
{
    \zR{宝钗道：“那可使不得。}
}
{
    \eR{[English source not present in the checked-in Bencraft/Joly files.]}
}
{
}

\Verse{230}
{
    \zR{已经报了，怎么能息呢？”}
}
{
    \eR{[English source not present in the checked-in Bencraft/Joly files.]}
}
{
}

\Verse{231}
{
    \zR{周瑞家的等人大家做好做歹的劝说：“若要息事，除非 夏亲家太太自己出去拦验，我们不提长短罢了。”}
}
{
    \eR{[English source not present in the checked-in Bencraft/Joly files.]}
}
{
}

\Verse{232}
{
    \zR{贾琏在外也将他儿子吓住。}
}
{
    \eR{[English source not present in the checked-in Bencraft/Joly files.]}
}
{
}

\Verse{233}
{
    \zR{他情 愿迎到刑部具结拦验，众人依允。}
}
{
    \eR{[English source not present in the checked-in Bencraft/Joly files.]}
}
{
}

\Verse{234}
{
    \zR{薛姨妈命人买棺成殓，不提。}
}
{
    \eR{[English source not present in the checked-in Bencraft/Joly files.]}
}
{
}

\Verse{235}
{
    \zR{且说贾雨村升了京兆府尹，兼管税务。}
}
{
    \eR{[English source not present in the checked-in Bencraft/Joly files.]}
}
{
}

\Verse{236}
{
    \zR{一日，出都查勘开垦地亩，路过知机县， 到了急流津，正要渡过彼岸，因待人夫，暂且停轿。}
}
{
    \eR{[English source not present in the checked-in Bencraft/Joly files.]}
}
{
}

\Verse{237}
{
    \zR{只见村旁有一座小庙，墙壁坍 颓，露出几株古松，倒也苍老。}
}
{
    \eR{[English source not present in the checked-in Bencraft/Joly files.]}
}
{
}

\Verse{238}
{
    \zR{雨村下轿，闲步进庙，但见庙内神像，金身脱落， 殿宇歪斜，旁有断碣，字迹模糊，也看不明白。}
}
{
    \eR{[English source not present in the checked-in Bencraft/Joly files.]}
}
{
}

\Verse{239}
{
    \zR{意欲行至后殿，只见一株翠柏下荫 着一间茅庐，庐中有一个道士，合眼打坐。}
}
{
    \eR{[English source not present in the checked-in Bencraft/Joly files.]}
}
{
}

\Verse{240}
{
    \zR{雨村走近看时，面貌甚熟，想着倒像在 那里见过的，一时再想不起来。}
}
{
    \eR{[English source not present in the checked-in Bencraft/Joly files.]}
}
{
}

\Verse{241}
{
    \zR{从人便欲吆喝，雨村止住，徐步向前，叫一声“老 道”。}
}
{
    \eR{[English source not present in the checked-in Bencraft/Joly files.]}
}
{
}

\Verse{242}
{
    \zR{那道士双眼略启，微微的笑道：“贵官何事？”}
}
{
    \eR{[English source not present in the checked-in Bencraft/Joly files.]}
}
{
}

\Verse{243}
{
    \zR{雨村便道：“本府出都查勘事 件，路过此地，见老道静修自得，想来道行深通，意欲冒昧请教。”}
}
{
    \eR{[English source not present in the checked-in Bencraft/Joly files.]}
}
{
}

\Verse{244}
{
    \zR{那道人说：“来 自有地，去自有方。”}
}
{
    \eR{[English source not present in the checked-in Bencraft/Joly files.]}
}
{
}

\Verse{245}
{
    \zR{雨村知是有些来历的，便长揖请问：“老道从何处焚修，在此 结庐？}
}
{
    \eR{[English source not present in the checked-in Bencraft/Joly files.]}
}
{
}

\Verse{246}
{
    \zR{此庙何名？}
}
{
    \eR{[English source not present in the checked-in Bencraft/Joly files.]}
}
{
}

\Verse{247}
{
    \zR{庙中共有几人？}
}
{
    \eR{[English source not present in the checked-in Bencraft/Joly files.]}
}
{
}

\Verse{248}
{
    \zR{或欲真修，岂无名山？}
}
{
    \eR{[English source not present in the checked-in Bencraft/Joly files.]}
}
{
}

\Verse{249}
{
    \zR{或欲结缘，何不通衢？”}
}
{
    \eR{[English source not present in the checked-in Bencraft/Joly files.]}
}
{
}

\Verse{250}
{
    \zR{那道 人道：“‘葫芦’尚可安身，何必名山结舍？}
}
{
    \eR{[English source not present in the checked-in Bencraft/Joly files.]}
}
{
}

\Verse{251}
{
    \zR{庙名久隐，断碣犹存，形影相随，何须 修募？}
}
{
    \eR{[English source not present in the checked-in Bencraft/Joly files.]}
}
{
}

\Verse{252}
{
    \zR{岂似那‘玉在椟中求善价，钗于匣内待时飞’之辈耶！”}
}
{
    \eR{[English source not present in the checked-in Bencraft/Joly files.]}
}
{
}

\Verse{253}
{
    \zR{雨村原是个颖悟人， 初听见“葫芦”两字，后闻“钗玉”一对，忽然想起甄士隐的事来，重复将那道士
        端详一回，见他容貌依然，便屏退从人，问道：“君家莫非甄老先生么？”}
}
{
    \eR{[English source not present in the checked-in Bencraft/Joly files.]}
}
{
}

\Verse{254}
{
    \zR{那道人 微微笑道：“什么‘真’？}
}
{
    \eR{[English source not present in the checked-in Bencraft/Joly files.]}
}
{
}

\Verse{255}
{
    \zR{什么‘假’？}
}
{
    \eR{[English source not present in the checked-in Bencraft/Joly files.]}
}
{
}

\Verse{256}
{
    \zR{要知道‘真’即是‘假’，‘假’即是‘真’。”}
}
{
    \eR{[English source not present in the checked-in Bencraft/Joly files.]}
}
{
}

\Verse{257}
{
    \zR{雨村听说出“贾”字来，益发无疑，便从新施礼，道：“学生自蒙慨赠到都，托庇 获隽公车，受任贵乡，始知老先生超悟尘凡，飘举仙境。}
}
{
    \eR{[English source not present in the checked-in Bencraft/Joly files.]}
}
{
}

\Verse{258}
{
    \zR{学生虽溯洄思切，自念风 尘俗吏，末由再睹仙颜，今何幸于此处相遇!求老仙翁指示愚蒙。}
}
{
    \eR{[English source not present in the checked-in Bencraft/Joly files.]}
}
{
}

\Verse{259}
{
    \zR{倘荷不弃，京寓 甚近，学生当得供奉，得以朝夕聆教。”}
}
{
    \eR{[English source not present in the checked-in Bencraft/Joly files.]}
}
{
}

\Verse{260}
{
    \zR{那道人也站起来回礼，道：“我于蒲团之外， 不知天地间尚有何物。}
}
{
    \eR{[English source not present in the checked-in Bencraft/Joly files.]}
}
{
}

\Verse{261}
{
    \zR{适才尊官所言，贫道一概不解。”}
}
{
    \eR{[English source not present in the checked-in Bencraft/Joly files.]}
}
{
}

\Verse{262}
{
    \zR{说毕依旧坐下。}
}
{
    \eR{[English source not present in the checked-in Bencraft/Joly files.]}
}
{
}

\Verse{263}
{
    \zR{雨村复又 心疑：“想去若非士隐，何貌言相似若此？}
}
{
    \eR{[English source not present in the checked-in Bencraft/Joly files.]}
}
{
}

\Verse{264}
{
    \zR{离别来十九载，面色如旧，必是修炼有成， 未肯将前身说破。}
}
{
    \eR{[English source not present in the checked-in Bencraft/Joly files.]}
}
{
}

\Verse{265}
{
    \zR{但我既遇恩公，又不可当面错过。}
}
{
    \eR{[English source not present in the checked-in Bencraft/Joly files.]}
}
{
}

\Verse{266}
{
    \zR{看来不能以富贵动之，那妻女 之私更不必说了。”}
}
{
    \eR{[English source not present in the checked-in Bencraft/Joly files.]}
}
{
}

\Verse{267}
{
    \zR{想罢，又道：“仙师既不肯说破前因，弟子于心何忍！”}
}
{
    \eR{[English source not present in the checked-in Bencraft/Joly files.]}
}
{
}

\Verse{268}
{
    \zR{正要下 礼，只见从人进来禀说：“天色将晚，快请渡河。”}
}
{
    \eR{[English source not present in the checked-in Bencraft/Joly files.]}
}
{
}

\Verse{269}
{
    \zR{雨村正无主意，那道人道：“请 尊官速登彼岸，见面有期，迟则风浪顿起。}
}
{
    \eR{[English source not present in the checked-in Bencraft/Joly files.]}
}
{
}

\Verse{270}
{
    \zR{果蒙不弃，贫道他日尚在渡头候教。”}
}
{
    \eR{[English source not present in the checked-in Bencraft/Joly files.]}
}
{
}

\Verse{271}
{
    \zR{说毕，仍合眼打坐。}
}
{
    \eR{[English source not present in the checked-in Bencraft/Joly files.]}
}
{
}

\Verse{272}
{
    \zR{雨村无奈，只得辞了道人出庙。}
}
{
    \eR{[English source not present in the checked-in Bencraft/Joly files.]}
}
{
}

\Verse{273}
{
    \zR{正要过渡，只见一人飞奔而来。}
}
{
    \eR{[English source not present in the checked-in Bencraft/Joly files.]}
}
{
}

\Verse{274}
{
    \zR{未知何人，下回分解。}
}
{
    \eR{[English source not present in the checked-in Bencraft/Joly files.]}
}
{
}

\Verse{275}
{
    \zR{(本章完)}
}
{
    \eR{[English source not present in the checked-in Bencraft/Joly files.]}
}
{
}
